\subsection{Induzione e somme finite}
\label{sec:esercizi-fondamenti-induzione}

\begin{exercise}[Somma di una progressione]
\label{ex:fondamenti-somma-progressione}
\difficoltaesercizio{1}
Dimostrare che, per ogni $n\ge1$,
\[
  \sum_{k=1}^{n}(2k+1)=n(n+2).
\]
\rimandoteoria{thm:principio-induzione}
\end{exercise}
\begin{proof}[Soluzione dell'esercizio]
Si pone
\[
  A:=\left\{n\in\N:n\ge1\ \text{e}\
  \sum_{k=1}^{n}(2k+1)=n(n+2)\right\}.
\]

\emph{Passo base.} Per $n=1$,
\[
  \sum_{k=1}^{1}(2k+1)=3=1\cdot(1+2),
\]
quindi $1\in A$.

\emph{Passo induttivo.} Sia $n\ge1$ e si supponga $n\in A$. Allora
\[
\begin{aligned}
  \sum_{k=1}^{n+1}(2k+1)
    &=\sum_{k=1}^{n}(2k+1)+\bigl(2(n+1)+1\bigr)\\
    &=n(n+2)+(2n+3)\\
    &=(n+1)(n+3).
\end{aligned}
\]
Pertanto $n\in A\Rightarrow n+1\in A$ e la tesi segue per induzione.
\end{proof}

\begin{exercise}[Somma di prodotti consecutivi]
\label{ex:fondamenti-somma-prodotti-consecutivi}
\difficoltaesercizio{1}
Dimostrare che, per ogni $n\ge1$,
\[
  \sum_{k=1}^{n}k(k+1)=\frac{n(n+1)(n+2)}{3}.
\]
\rimandoteoria{thm:principio-induzione}
\end{exercise}
\begin{proof}[Soluzione dell'esercizio]
Si pone
\[
  A:=\left\{n\in\N:n\ge1\ \text{e}\
    \sum_{k=1}^{n}k(k+1)=\frac{n(n+1)(n+2)}{3}\right\}.
\]

\emph{Passo base.} Per $n=1$,
\[
  \sum_{k=1}^{1}k(k+1)=2=\frac{1\cdot2\cdot3}{3},
\]
quindi $1\in A$.

\emph{Passo induttivo.} Sia $n\ge1$ e si supponga $n\in A$. Allora
\[
\begin{aligned}
  \sum_{k=1}^{n+1}k(k+1)
    &=\sum_{k=1}^{n}k(k+1)+(n+1)(n+2)\\
    &=\frac{n(n+1)(n+2)}{3}+(n+1)(n+2)\\
    &=\frac{(n+1)(n+2)(n+3)}{3}.
\end{aligned}
\]
Pertanto $n\in A\Rightarrow n+1\in A$ e la tesi segue per induzione.
\end{proof}

\begin{exercise}[Somma geometrica con ragione fissata]
\label{ex:fondamenti-somma-geometrica-ragione-tre}
\difficoltaesercizio{1}
Dimostrare, indicando esplicitamente il passo base e il passo induttivo,
che, per ogni $n\in\N$,
\[
  \sum_{k=0}^{n}3^k=\frac{3^{n+1}-1}{2}.
\]
\rimandoteoria{thm:principio-induzione}
\end{exercise}
\begin{proof}[Soluzione dell'esercizio]
Si pone
\[
  A:=\left\{n\in\N:\sum_{k=0}^{n}3^k=\frac{3^{n+1}-1}{2}\right\}.
\]

\emph{Passo base.} Per $n=0$,
\[
  \sum_{k=0}^{0}3^k=1=\frac{3-1}{2},
\]
quindi $0\in A$.

\emph{Passo induttivo.} Sia $n\in\N$ e si supponga $n\in A$. Allora
\[
\begin{aligned}
  \sum_{k=0}^{n+1}3^k
    &=\sum_{k=0}^{n}3^k+3^{n+1}\\
    &=\frac{3^{n+1}-1}{2}+3^{n+1}\\
    &=\frac{3^{n+2}-1}{2}.
\end{aligned}
\]
Pertanto $n\in A\Rightarrow n+1\in A$ e la tesi segue per induzione.
\end{proof}

\clearpage

\begin{exercise}[Contronominale e parità dei cubi]
\label{ex:fondamenti-parita-del-cubo}
\difficoltaesercizio{1}
Dimostrare per contronominale che, per ogni $m\in\Z$,
\[
  m^3\text{ pari}\quad\Longrightarrow\quad m\text{ pari}.
\]
Utilizzare la rappresentazione $m=2k+1$ per gli interi dispari.
\rimandoteoria{sec:lez01-insiemi-dimostrazioni}
\end{exercise}
\begin{proof}[Soluzione dell'esercizio]
Si dimostra la contronominale. Sia $m$ dispari, ossia
$m=2k+1$ per un $k\in\Z$. Allora
\[
\begin{aligned}
  m^3&=(2k+1)^3\\
     &=8k^3+12k^2+6k+1\\
     &=2(4k^3+6k^2+3k)+1.
\end{aligned}
\]
Poiché $4k^3+6k^2+3k\in\Z$, anche $m^3$ è dispari.
La tesi segue per contronominale.
\end{proof}

\begin{exercise}[Confronto tra una potenza e l'indice]
\label{ex:fondamenti-potenza-indice}
\difficoltaesercizio{1}
Dimostrare per induzione che $2^n\ge n+1$ per ogni $n\in\N$.
Stabilire per quali $n$ si ha uguaglianza.
\rimandoteoria{thm:principio-induzione}
\end{exercise}
\begin{proof}[Soluzione dell'esercizio]
Si pone
\[
  A:=\{n\in\N:2^n\ge n+1\}.
\]

\emph{Passo base.} Per $n=0$,
\[
  2^0=1=0+1,
\]
quindi $0\in A$.

\emph{Passo induttivo.} Sia $n\in\N$ e si supponga $n\in A$. Allora
\[
  2^{n+1}=2\cdot2^n\ge2(n+1)\ge n+2.
\]
Pertanto $n\in A\Rightarrow n+1\in A$ e la tesi segue per induzione.

L'uguaglianza vale per $n=0$ e $n=1$; per $n\ge2$ la disuguaglianza
è stretta.
\end{proof}

\clearpage

\begin{exercise}[Confronto tra una potenza e un polinomio]
\label{ex:fondamenti-potenza-polinomio}
\difficoltaesercizio{1}
Dimostrare per induzione che, per ogni $n\in\N$,
\[
  3^n\ge2n^2+1.
\]
\rimandoteoria{thm:principio-induzione}
\end{exercise}
\begin{proof}[Soluzione dell'esercizio]
Si pone
\[
  A:=\{n\in\N:3^n\ge2n^2+1\}.
\]

\emph{Passo base.} Per $n=0$,
\[
  3^0=1=2\cdot0^2+1,
\]
quindi $0\in A$.

\emph{Passo induttivo.} Sia $n\in\N$ e si supponga $n\in A$. Allora
\[
  3^{n+1}=3\cdot3^n\ge6n^2+3\ge2(n+1)^2+1,
\]
poiché $4n(n-1)\ge0$ per ogni $n\in\N$.
Pertanto $n\in A\Rightarrow n+1\in A$ e la tesi segue per induzione.
\end{proof}

\begin{exercise}[Somma dei cubi]
\label{ex:fondamenti-somma-dei-cubi}
\difficoltaesercizio{1}
Dimostrare mediante induzione che, per ogni $n\ge1$,
\[
  \sum_{k=1}^{n}k^3=\left(\frac{n(n+1)}{2}\right)^2.
\]
Esplicitare la fattorizzazione necessaria nel passo induttivo.
\rimandoteoria{thm:principio-induzione}
\end{exercise}
\begin{proof}[Soluzione dell'esercizio]
Si pone
\[
  A:=\left\{n\in\N:n\ge1\ \text{e}\
    \sum_{k=1}^{n}k^3=\left(\frac{n(n+1)}{2}\right)^2\right\}.
\]

\emph{Passo base.} Per $n=1$,
\[
  \sum_{k=1}^{1}k^3=1=\left(\frac{1\cdot2}{2}\right)^2,
\]
quindi $1\in A$.

\emph{Passo induttivo.} Sia $n\ge1$ e si supponga $n\in A$. Allora
\[
\begin{aligned}
  \sum_{k=1}^{n+1}k^3
    &=\sum_{k=1}^{n}k^3+(n+1)^3\\
    &=\left(\frac{n(n+1)}{2}\right)^2+(n+1)^3\\
    &=\frac{(n+1)^2\bigl(n^2+4(n+1)\bigr)}{4}\\
    &=\left(\frac{(n+1)(n+2)}{2}\right)^2.
\end{aligned}
\]
Pertanto $n\in A\Rightarrow n+1\in A$ e la tesi segue per induzione.
\end{proof}
